\section{Execution Model}
\label{sec:execution-model}

The memory model says what the store contains. The execution model says how
\glspl{process} walk along the timeline reading and writing it, and what each
must do so its neighbours can make progress.

% ===========================================================================
\subsection{Processes and the shape of a program}
\label{sec:exec-shape}

A DSTM program is one sequential program executed by every process. Each runs
a single thread of control, occupies one point on the timeline at a time, and
hosts no virtualised sub-entities: no logical processes, no migratable agents.
Agents, if wanted, are built above DSTM with their state in the \gls{store}.

The canonical shape is a loop over ticks:

\begin{lstlisting}
for (dstm_time_t t = 0; t < T; t++) {
        /* 1. consume: read what neighbours made visible at t */
        /* 2. compute: local work, arbitrary duration                */
        /* 3. produce: publish results, visible at t + delta         */
}
dstm_barrier(last_time);
\end{lstlisting}

There is no \gls{barrier} at the tick boundary. \Glspl{process} leave the loop
body at different wall-clock instants and reach tick $t+1$ at different
instants; nothing resynchronises them.

\begin{definition}[Program]\label{def:program}
The \emph{program} of rank $p$ is the sequence of \glspl{operation} $p$ issues,
ordered by $\po$ (\cref{def:history}). A program is \emph{well clocked} when
that sequence has nondecreasing \glspl{origin-time}.
\end{definition}

The runtime does not check well-clockedness \devmark{dev:monotone-origin}. It
is a program restriction (\cref{res:monotone}), and essentially every
guarantee here is conditional on it.

% ===========================================================================
\subsection{Clocks and their publication}
\label{sec:exec-clocks}

Position on the \gls{logical-time} timeline is communicated implicitly: there
is no ``advance to time $t$'' call. A process declares where it is by issuing
an operation with origin time $t$; when that exceeds the clock, the runtime
advances it and broadcasts a heartbeat to the \gls{neighbourhood}.

\begin{definition}[Clock advance]\label{def:clock-advance}
Issuing an operation with origin time $t > \clk{p}$ has the effect
\[
  \clk{p} \mathrel{{:}{=}} t,
  \qquad
  \text{heartbeat } \langle p, t \rangle \text{ sent to every } q \in \nbr{p}.
\]
Issuing with $t \le \clk{p}$ leaves the clock unchanged and sends nothing.
\end{definition}

Two things follow.

First, \textbf{a silent process blocks its neighbours.} A process that stops issuing
operations stops advancing its clock, so its neighbours' \glspl{known-time} for it
stop rising, so the \glspl{version} those neighbours hold cannot be published. A process
with nothing to say must still say it, by continuing to issue at increasing
origin times or by entering the \gls{termination} barrier.

Second, \textbf{the clock tracks origin time, not \gls{effective-time}.} A put with
origin $t$ and \gls{delta} $\delta$ advances the clock to $t$, not to $t+\delta$
\devmark{dev:clock-conflation}. This is the correct choice --- the clock's job
is to report where the \emph{issuer} is, and the issuer is at $t$ --- but it
means the largest effective time a process has produced can exceed its clock by
an unbounded amount, which is exactly what \cref{res:horizon} constrains.

\begin{definition}[Known-time update]\label{def:known-update}
Every message carries its sender's identity $q$ and a stamp $t \le \clk{q}$.
On receipt at $p$: $\kn{p}{q} \mathrel{{:}{=}} \max(\kn{p}{q}, t)$.
\end{definition}

The stamp is a \emph{lower} bound on the sender's clock, not necessarily equal
to it: a reply to a read is stamped with the read's \gls{effective-time}
rather than the responder's clock. Taking the maximum makes known times
monotone, and \cref{inv:clock-soundness} makes them sound. Underestimating is
harmless where overestimating would not be:

\begin{proposition}[Conservatism of known times]\label{prop:conservative}
Both promotion guards of \cref{sec:publication} are monotone in the known
times: lowering any $\kn{p}{q}$ can only prevent a promotion, never cause one.
Consequently an implementation that underestimates neighbour clocks is
\emph{safe but possibly slow}, whereas one that overestimates them is
\emph{unsound}.
\end{proposition}

% ===========================================================================
\subsection{Issue, completion and blocking}
\label{sec:exec-blocking}

The two \glspl{operation} differ in the gap between \gls{issue} and
\gls{completion}, and the asymmetry is deliberate.

\begin{description}[leftmargin=0.9cm,style=nextline,font=\normalfont\bfseries]
\item[A put is asynchronous.]
  It returns as soon as the payload has been copied out of the caller's
  buffer. It carries no result and cannot report a delivery failure. Its
  effect on the store is settled later, possibly long after the issuing process
  has moved on. A program therefore has no way to observe that its own put has
  landed, other than by reaching \gls{quiescence}.

\item[A get is synchronous.]
  It does not return until the requested version has been published and its
  bytes copied into the caller's buffer. While waiting it drives the runtime's
  progress engine, so a blocked process continues to serve its neighbours'
  requests, forward messages and promote its own versions. This is what keeps
  a blocked get from becoming a deadlock in the common case.
\end{description}

\begin{remark}[Blocking is where drift is absorbed]
A get is the only place a DSTM program waits, and its duration is exactly how
far the reader has run ahead of the producer it reads. Versioning, lag
bounds and \glspl{heartbeat} all exist to make that wait as short as the
application's real data dependence requires.
\end{remark}

% ===========================================================================
\subsection{The task grid and drift}
\label{sec:exec-grid}

Picture an execution as a grid: one row per rank, one column per tick, one
\emph{task} per \gls{cell}. The task $\task{p}{t}$ is the work $\proc{p}$
performs at origin time $t$ --- one loop iteration. \Glspl{operation} create
edges between tasks, weighted by their \gls{delta}.

\begin{definition}[Task]\label{def:task}
A \emph{task} $\task{p}{t}$ is the set of operations issued by $p$ with origin
time $t$, together with the local computation between them. Its \emph{work}
$\work{p}{t} \ge 0$ is the wall-clock duration of that computation.
\end{definition}

\begin{figure}[htbp]
\centering
\begin{tikzpicture}[
  x=17mm, y=-12mm,
  every node/.style={font=\scriptsize},
  tk/.style={draw,rounded corners=1.5pt,minimum width=11mm,minimum height=6mm,
             fill=resbg,draw=dstmblue,inner sep=1pt},
  >={Stealth[round]},
]
  % tick axis
  \foreach \t in {0,1,2,3,4} {
    \node[text=dstmgrey] at (\t,-0.75) {$t=\t$};
  }
  % process rows
  \foreach \r/\name/\last in {0/r_0/4, 1/r_1/3, 2/r_2/2, 3/r_3/3} {
    \node[text=dstmgrey,anchor=east] at (-0.55,\r) {$\name$};
    \draw[dstmgrey!40] (-0.35,\r) -- (4.6,\r);
    \foreach \t in {0,...,\last} { \node[tk] (n\r-\t) at (\t,\r) {$\t$}; }
  }
  % put edges: forward in time, to a neighbour  (solid)
  \draw[->,thick,dstmblue] (n0-1) -- node[above right,inner sep=1pt]{$+1$} (n1-2);
  \draw[->,thick,dstmblue] (n1-1) -- node[above right,inner sep=1pt]{$+1$} (n2-2);
  \draw[->,thick,dstmblue] (n3-1) -- node[below right,inner sep=1pt]{$+1$} (n2-2);
  \draw[->,thick,dstmblue] (n0-2) -- (n3-3);
  % get edges: backward in time, from a neighbour (dashed)
  \draw[->,thick,dashed,dstmred] (n1-2) to[bend left=18] (n0-3);
  \draw[->,thick,dashed,dstmred] (n2-2) to[bend right=18] (n3-3);
  % drift bracket
  \draw[decorate,decoration={brace,amplitude=4pt},dstmamber,thick]
       (2.55,2.35) -- (4.45,2.35)
       node[midway,below=4pt,text=dstmamber]{drift of $r_0$ over $r_2$};
\end{tikzpicture}
\caption{An execution as a task grid, for a four-process ring with all lags
equal to $1$. Solid edges are puts (forward in time, into a neighbour's
region); dashed edges are gets (backward in time, out of a neighbour's
region). Process $r_0$ has reached tick $4$ while $r_2$ is still at tick $2$:
they have drifted apart by two ticks and neither has rolled back or waited at
a barrier. The edges, not the rows, are what constrain the schedule.}
\label{fig:task-grid}
\end{figure}

\Cref{fig:task-grid} makes the central claim visible: rows are independent,
and the only coupling is the edges, every one of which spans at least one tick
because every \gls{lag} is at least one. A process is never blocked by
\emph{where} its neighbours are, only by \emph{what} they have already
produced.

Drift is not bounded by the model, but it is bounded by the
\gls{comm-graph} whenever the program communicates each tick. That hypothesis
is worth naming.

\begin{definition}[Coupled schedule]\label{def:coupled}
A schedule is \emph{coupled} when, at every tick $t$ it executes, each process
$\proc{p}$ completes a \gls{get} whose target \gls{version} cannot be
published until $\clk{q} \ge t - \lag{p}{q}$ for every $q \in \nbr{p}$. Write
$T_p$ for the greatest tick $\proc{p}$ has completed.
\end{definition}

Both standard idioms are coupled, for slightly different reasons. Under the
\gls{push-model} a process reads its own region at $\delta = 0$, so publishing
$\ver{p}{t}$ requires $\kn{p}{q} + \lag{p}{q} > t$ directly. Under the
\gls{pull-model} it reads $\ver{q}{t - \lag{p}{q}}$, whose publication
requires $\clk{q} \ge t - \lag{p}{q}$ by the local-clock clause.
\Cref{prop:drift-bound} turns this into a bound computed from $\cgraph$ alone,
and \cref{cor:space-drift} turns that into a bound on memory.

% ===========================================================================
\subsection{Session lifecycle}
\label{sec:exec-lifecycle}

\begin{figure}[htbp]
\centering
\resizebox{\textwidth}{!}{%
\begin{tikzpicture}[
  node distance=6mm and 9mm,
  every node/.style={font=\scriptsize},
  ph/.style={draw,rounded corners=2pt,minimum height=8mm,align=center,
             inner xsep=4pt,draw=dstmblue,fill=resbg},
  >={Stealth[round]},
]
  \node[ph] (init)  {\code{dstm\_init}};
  \node[ph,right=of init]  (start) {\code{dstm\_start}$(\Dmat)$};
  \node[ph,right=of start] (mk)    {\code{darray/dgraph}\\\code{\_create}};
  \node[ph,right=of mk]    (loop)  {tick loop\\\code{put}/\code{get}};
  \node[ph,right=of loop]  (fini)  {\code{dstm\_fini}$(\ell)$};
  \node[ph,right=of fini]  (final) {\code{dstm\_finalize}};
  \foreach \a/\b in {init/start,start/mk,mk/loop,loop/fini,fini/final}
    { \draw[->,thick] (\a) -- (\b); }
  \draw[->,thick,dstmteal] (fini.south) .. controls +(0,-7mm) and +(0,-7mm) ..
       node[below,text=dstmteal]{further sessions} (start.south);

  \begin{scope}[on background layer]
    \node[fit=(start)(fini),draw=dstmteal,dashed,rounded corners,
          inner sep=5pt,label={[text=dstmteal,font=\scriptsize]above:session}] {};
  \end{scope}
\end{tikzpicture}}
\caption{Session lifecycle. The \gls{comm-graph} and every region are fixed
between \code{dstm\_start} and \code{dstm\_fini}; changing topology means
closing the session and opening another.}
\label{fig:lifecycle}
\end{figure}

The restrictions attach to different scopes: \cref{res:collective} to the
whole process lifetime; \crefrange{res:monotone}{res:phase} and
\cref{res:write-disjoint} to a single \gls{session}; \cref{res:terminate} to
the transition out of one.

% ===========================================================================
\subsection{The two synchronisation idioms}
\label{sec:exec-idioms}

The API admits many \gls{delta} choices; applications use two, and both
benchmark applications in \cite{dstm2026} use one of them exclusively.

\begin{center}
\small
\begin{tabular}{@{}lll@{}}
\toprule
 & \textbf{Pull} & \textbf{Push} \\
\midrule
Producer $p$ writes & own region, $\delta = 0$
                    & neighbour $q$'s region, $\delta = +\lag{p}{q}$ \\
Consumer $q$ reads  & $p$'s region, $\delta = -\lag{p}{q}$
                    & own region, $\delta = 0$ \\
Data movement       & on demand, at read time & eager, at write time \\
Messages per pair   & two (request, reply)    & one \\
\bottomrule
\end{tabular}
\end{center}

They look like opposites, but the value produced by $\proc{p}$ at tick $t$
becomes available to $\proc{q}$ at tick $t + \lag{p}{q}$ in both cases: the
\gls{lag}, not the direction of data movement, sets the pipeline depth.
\Cref{thm:pipeline-depth} states this precisely. The choice is therefore a
communication-cost decision --- one message versus two --- never a semantic
one.